\section{Desarrollo}\label{sec:desarrollo}

\subsection{Análisis Teórico de Algoritmos de Ordenamiento}\label{sec:ordenamiento_dv}
En esta sección se detallan los algoritmos de ordenamiento implementados bajo el paradigma de \textit{Divide y Vencerás}. Se analizará su funcionamiento de manera teórica, desglosando sus pseudocódigos y estudiando su complejidad asintótica global.

\subsubsection{Merge Sort: Clásico y Optimizado}\label{sec:merge_sort}
El algoritmo Merge Sort basa su funcionamiento en dividir el arreglo por la mitad repetidamente hasta obtener arreglos de un solo elemento, para luego fusionarlos (merge) de manera ordenada.
\vspace{0.3cm}
\par \noindent\textit{[Aquí se insertará el pseudocódigo, la explicación teórica y el análisis de la optimización del umbral combinando con Insertion Sort].}

\subsubsection{Quick Sort y sus Variantes de Pivote}\label{sec:quick_sort}
Quick Sort selecciona un elemento como pivote y particiona el arreglo (en este caso implementado según el esquema de partición de Lomuto), ubicando los elementos menores a su izquierda y los mayores a su derecha, para luego procesar los sub-arreglos recursivamente.
\vspace{0.3cm}
\par \noindent\textit{[Aquí se insertará el pseudocódigo, la explicación de las 4 heurísticas de selección de pivote (Último, Primero, Aleatorio, Mediana de tres) y el análisis de su impacto en el mejor, peor y caso promedio].}


\subsection{Análisis Teórico de Algoritmos de Búsqueda}\label{sec:busqueda_dv}
A los algoritmos clásicos de la primera iteración del proyecto se ha sumado un conjunto de algoritmos de búsqueda basados en reducir progresiva o logarítmicamente el espacio restante.

\subsubsection{Búsqueda Binaria: Recursiva y por Rango}\label{sec:busqueda_binaria_rango}
Basada en evaluar permanentemente el valor medio de un sector de estudio asumiendo arreglos ya ordenados, la versión recursiva delega el ciclo a la pila del sistema. Su extensión natural permite no solo atinar a un elemento sino determinar el límite inferior y superior de múltiples coincidencias de este en el arreglo.
\vspace{0.3cm}
\par \noindent\textit{[Aquí se insertará el pseudocódigo, análisis y consideraciones teóricas correspondientes].}

\subsubsection{Búsqueda Exponencial}\label{sec:busqueda_exponencial}
La búsqueda exponencial halla un rango propicio delimitado escalando en potencias de 2, para luego ceder y aplicar una búsqueda binaria local en el tramo hallado.
\vspace{0.3cm}
\par \noindent\textit{[Aquí se insertará el pseudocódigo y la demostración en complejidad analítica].}

\subsubsection{Búsqueda por Interpolación}\label{sec:busqueda_interpolacion}
Mimetizando la forma en que los seres humanos buscan en un diccionario, se realiza una aproximación probabilística del punto objetivo a partir de los extremos del área evaluada para dar atajos rápidos. 
\vspace{0.3cm}
\par \noindent\textit{[Aquí se insertará el pseudocódigo y la complejidad asintótica, profundizando en por qué su peor caso depende directamente de la variabilidad del conjunto].}


\subsection{Análisis Teórico de Algoritmos de Selección}\label{sec:seleccion_dv}

\subsubsection{Quick Select}\label{sec:quick_select}
Basado fuertemente en el esquema de Quick Sort, se detiene prematuramente en la recursividad, desechando la porción del arreglo donde matemáticamente no puede existir el \textit{k-ésimo} elemento buscado.
\vspace{0.3cm}
\par \noindent\textit{[Aquí se insertará el pseudocódigo de manera formal y el análisis de la consecuente complejidad en el mejor y peor de los casos].}


\subsection{Resultados Experimentales}\label{sec:resultados_experimentales}
Esta sección contrasta los hallazgos teóricos de la complejidad con simulaciones obtenidas de los datasets estructurados.

\subsubsection{Impacto del Umbral en Merge Sort}\label{sec:exp_merge_sort_umbral}
\vspace{0.3cm}
\par \noindent\textit{[Aquí se analizarán gráficamente con plots comparativos los distintos umbrales para dirimir cómo rinde frente a la versión clásica pura].}

\subsubsection{Comparativa de Pivotes en Quick Sort}\label{sec:exp_quick_sort_pivotes}
\vspace{0.3cm}
\par \noindent\textit{[Aquí se someterán a prueba, documentando en gráficas, las metodologías de pivotaje estudiadas frente al peor, mejor y caso promedio de los datos de los usuarios].}

\subsubsection{Desempeño en Algoritmos de Búsqueda}\label{sec:exp_busquedas}
\vspace{0.3cm}
\par \noindent\textit{[Aquí irán las observaciones empíricas relacionadas con la agilidad y las caídas bruscas causadas por peores casos medibles].}

\subsubsection{Contraste: Divide y Vencerás frente a Incrementales (Tarea 1)}\label{sec:exp_vs_t1}
\vspace{0.3cm}
\par \noindent\textit{[Aquí se exhibirá en un marco gráfico comparativo la monumental diferencia del desempeño que hay entre O($n^2$) clásicos preexistentes frente a la agilidad del esquema de orden O($n \log n$)].}


\subsection{Aplicación Práctica y Funcionalidades}\label{sec:aplicacion_practica}
\vspace{0.3cm}
\par \noindent\textit{[Aquí se documentará de forma visual y teórica cómo la interfaz ha implementado las nuevas lógicas para listados de rango de puntajes, top atletas, ranking por N criterios o k-ésimo posicionado en el roster central del código fuente. Añadiendo una breve justificación sobre la selección del algoritmo para cada pantalla].}
